% from ega4-9-fr.tex
\oldpage[IV-3]{64}

\begin{corollary}[9.4.5]
\label{IV.9.4.5-fr}
Soient $f:X\to S$ un morphisme de présentation finie, $\mathcal F$,
$\mathcal G$ deux $\mathcal O_X$-Modules quasi-cohérents de présentation finie,
$u:\mathcal F\to\mathcal G$ un homomorphisme de $\mathcal O_X$-Modules.
Alors l'ensemble des $s\in S$ tels que $u_s$ soit injectif
\textnormal{(}resp. surjectif, bijectif, nul\textnormal{)} est localement
constructible.
\end{corollary}

Il suffit d'appliquer (9.4.4) aux suites
$0\to\mathcal F\xrightarrow{u}\mathcal G$,
$\mathcal F\xrightarrow{u}\mathcal G\to0$,
$0\to\mathcal F\xrightarrow{u}\mathcal G\to0$,
$\mathcal F\xrightarrow{u}\mathcal G\xrightarrow{1_\mathcal G}\mathcal G$.

\begin{corollary}[9.4.6]
\label{IV.9.4.6-fr}
Soient $f:X\to S$ un morphisme de présentation finie, $\mathcal F$ un
$\mathcal O_X$-Module quasi-cohérent de présentation finie. Soit $h$ une
section de $\mathcal F$ au-dessus de $X$~; pour tout $s\in S$, soit $h_s$ la
section correspondante de $\mathcal F_s$ au-dessus de $X_s$ \textnormal{(}pour
le morphisme projection $X_s\to X$\textnormal{)}. Alors, l'ensemble des
$s\in S$ tels que $h_s=0$ est localement constructible.
\end{corollary}

Il suffit de remarquer que $h$ correspond à un homomorphisme
$u:\mathcal O_X\to\mathcal F$ $(\mathbf 0_{\mathrm I},5.1.1)$ et $h_s$ à
l'homomorphisme $u_s$.

\begin{proposition}[9.4.7]
\label{IV.9.4.7-fr}
Soient $f:X\to S$ un morphisme de présentation finie, $\mathcal F$ un
$\mathcal O_X$-Module quasi-cohérent de présentation finie. L'ensemble $E$
\textnormal{(}resp. $E'$\textnormal{)} des $s\in S$ tels que $\mathcal F_s$
soit un $\mathcal O_{X_s}$-Module localement libre \textnormal{(}resp.
localement libre de rang $n$\textnormal{)} est localement constructible.
\end{proposition}

Si $X$ est un préschéma algébrique sur un corps $k$, $\mathcal F$ un
$\mathcal O_X$-Module cohérent, $k'$ une extension de $k$, alors, pour que
$\mathcal F$ soit localement libre (resp. localement libre de rang $n$), il
faut et il suffit qu'il en soit de même de $\mathcal F\otimes_k k'$, puisque la
projection $X_{(k')}\to X$ est un morphisme fidèlement plat (2.2.7). Autrement
dit, la condition (9.2.1, (i)) est vérifiée pour les propriétés dont on veut
démontrer la constructibilité, et il reste à vérifier (9.2.1, (ii))~; on peut
donc encore supposer que $S$ est affine, noethérien et intègre. Il y a de
nouveau quatre cas à envisager~:

$1^\circ$ $\eta\in E$. Il résulte de (8.5.5), appliqué suivant la méthode
générale de (8.1.2, \emph{a)}), qu'il existe un voisinage ouvert $U$ de $\eta$
dans $S$ tel que $\mathcal F|f^{-1}(U)$ soit localement libre~;
\emph{a fortiori} $\mathcal F_s$ est localement libre pour tout $s\in U$.

$2^\circ$ $\eta\in E'$. Même raisonnement que dans le $1^\circ$.

$3^\circ$ $\eta\in S-E$. Comme $\mathcal F_\eta$ est un
$\mathcal O_{X_\eta}$-Module cohérent, dire qu'il n'est pas localement libre
équivaut à dire qu'il \emph{n'est pas plat} sur $\mathcal O_{X_\eta}$
(Bourbaki, \emph{Alg. comm.}, chap.~II, \S~5, n\textsuperscript{o}~2,
cor.~2 du th.~1). Le fait que $S-E$ est un voisinage de $\eta$ résultera donc
du lemme plus général suivant (appliqué au cas où $g$ est l'identité)~:

\begin{lemma}[9.4.7.1]
\label{IV.9.4.7.1-fr}
Soient $S$ un préschéma intègre noethérien, de point générique $\eta$, $X$,
$Y$ deux $S$-préschémas de type fini sur $S$, $g:X\to Y$ un $S$-morphisme,
$\mathcal F$ un $\mathcal O_X$-Module cohérent. Si $\mathcal F_\eta$ n'est
pas $g_\eta$-plat, alors il existe un voisinage ouvert $U$ de $\eta$ dans $S$
tel que pour tout $s\in U$, $\mathcal F_s$ ne soit pas $g_s$-plat.
\end{lemma}

Compte tenu de (2.1.2) et de Bourbaki, \emph{Alg. comm.}, chap.~I\textsuperscript{er},
\S~2, n\textsuperscript{o}~3, Remarque~1, l'hypothèse signifie qu'il existe un
ouvert non vide $V$ de $Y_\eta$, et un homomorphisme injectif
$v:\mathcal M\to\mathcal N$ de $\mathcal O_V$-Modules \emph{cohérents}, tel
que l'homomorphisme
$1\otimes v:\mathcal F_\eta\otimes_{\mathcal O_V}\mathcal M\to
\mathcal F_\eta\otimes_{\mathcal O_V}\mathcal N$ ne soit pas injectif. On a
$V=Y_\eta\cap W$, où $W$ est ouvert dans $Y$ (\textbf{I}, 3.6.1), et il
résulte de (8.5.2, (i) et (ii)), appliqué suivant la méthode de
(8.1.2, \emph{a)}), qu'il existe un voisinage ouvert $U_0$ de $\eta$ dans $S$,
deux $\mathcal O_Z$-Modules

